%%%%%%%%%%%%%%%%%%%%%%%%%%%%%%%%%%%%%%%%%%%%%%%%%%%%%%%%%%%%%%%%%%%%%%%%%%%%%%%%
%% 
%% Abstract of your thesis in default document language
%% 

\cleardoubleoddpage%  Make sure to start abstract on a new odd page
\begin{abstract*}
This thesis evaluates the integration of \acf{xLSTM} architectures within the \acf{BIOT} framework for biosignal classification tasks. While \ac{BIOT}'s original architecture employs Linear Transformers for processing tokenized biosignal sequences, this work investigates whether \ac{xLSTM} variants can better capture temporal patterns in the data.

The study compares four model configurations on the \ac{TUEV} dataset: the original Linear Transformer-based \ac{BIOT}, \ac{mLSTM}, \ac{sLSTM}, and combined M+S-\ac{LSTM} variants. Experiments were conducted across multiple temporal window lengths (5s, 7s, 9s) using standard evaluation metrics including Balanced Accuracy, Cohen's Kappa, and Weighted F1 Score. The experimental setup maintained consistency with the original \ac{BIOT} evaluation protocol while incorporating architectural modifications to replace the Linear Transformer component with \ac{xLSTM} blocks.

Results show modest improvements for \ac{xLSTM} variants, with \ac{sLSTM} achieving the best performance of 52.26\% balanced accuracy on 7-second temporal windows, representing a 4.4\% relative improvement over the baseline Linear Transformer. However, statistical significance testing reveals that most performance differences are not statistically significant due to high inter-run variance. The \ac{mLSTM} variant demonstrates more stable performance across runs, while the combined architecture unexpectedly produces identical results to \ac{mLSTM} alone.

The findings indicate that while \ac{xLSTM} architectures can provide consistent numerical improvements in biosignal classification, the benefits are modest and limited by training variability. The study highlights the importance of statistical validation in architectural comparisons and reveals challenges in achieving robust performance gains in \ac{EEG}-based seizure classification tasks. This work contributes empirical evidence about the practical limitations and potential of modern sequence modeling architectures in clinical biosignal processing applications.
\end{abstract*}

%% 
%%%%%%%%%%%%%%%%%%%%%%%%%%%%%%%%%%%%%%%%%%%%%%%%%%%%%%%%%%%%%%%%%%%%%%%%%%%%%%%%


%%%%%%%%%%%%%%%%%%%%%%%%%%%%%%%%%%%%%%%%%%%%%%%%%%%%%%%%%%%%%%%%%%%%%%%%%%%%%%%%
%% 
%% Abstract of your thesis in an additional language
%% 

\cleardoubleoddpage%  Make sure to start abstract on a new odd page
\begin{otherlanguage}{ngerman}
\begin{abstract*}
Diese Arbeit evaluiert die Integration von \acf{xLSTM}-Architekturen innerhalb des \acf{BIOT}-Frameworks für die Klassifikation von Biosignaldaten. Während \ac{BIOT}s ursprüngliche Architektur Linear Transformer für die Verarbeitung tokenisierter Biosignal-Sequenzen verwendet, untersucht diese Arbeit, ob \ac{xLSTM}-Varianten zeitliche Muster in den Daten besser erfassen können.

Die Studie vergleicht vier Modellkonfigurationen auf dem \ac{TUEV}-Datensatz: das ursprüngliche Linear Transformer-basierte \ac{BIOT}, \ac{mLSTM}, \ac{sLSTM} und kombinierte M+S-\ac{LSTM}-Varianten. Experimente wurden über mehrere zeitliche Fensterlängen (5s, 7s, 9s) hinweg durchgeführt, wobei Standardevaluationsmetriken wie Balanced Accuracy, Cohen's Kappa und Weighted F1 Score verwendet wurden. Das experimentelle Setup behielt die Konsistenz mit dem ursprünglichen \ac{BIOT}-Evaluationsprotokoll bei, während architektonische Modifikationen eingeführt wurden, um die Linear Transformer-Komponente durch \ac{xLSTM}-Blöcke zu ersetzen.

Die Ergebnisse zeigen moderate Verbesserungen für \ac{xLSTM}-Varianten, wobei \ac{sLSTM} die beste Leistung von 52,26\% Balanced Accuracy bei 7-Sekunden-Zeitfenstern erreicht, was eine relative Verbesserung von 4,4\% gegenüber dem Baseline Linear Transformer darstellt. Statistische Signifikanztests zeigen jedoch, dass die meisten Leistungsunterschiede aufgrund hoher Varianz zwischen den Durchläufen nicht statistisch signifikant sind. Die \ac{mLSTM}-Variante zeigt stabilere Leistung über verschiedene Durchläufe hinweg, während die kombinierte Architektur unerwarteterweise identische Ergebnisse wie \ac{mLSTM} allein erzielt.

Die Befunde zeigen, dass \ac{xLSTM}-Architekturen zwar konsistente numerische Verbesserungen in der Biosignalklassifikation bieten können, die Vorteile jedoch bescheiden und durch Trainingsvariabilität begrenzt sind. Die Studie unterstreicht die Wichtigkeit statistischer Validierung bei Architekturvergleichen und zeigt Herausforderungen beim Erreichen robuster Leistungsverbesserungen in \ac{EEG}-basierten Anfallsklassifikationsaufgaben auf. Diese Arbeit trägt empirische Evidenz über die praktischen Grenzen und das Potenzial moderner Sequenzmodellierungsarchitekturen in klinischen Biosignalverarbeitungsanwendungen bei.
\end{abstract*}
\end{otherlanguage}

%% 
%%%%%%%%%%%%%%%%%%%%%%%%%%%%%%%%%%%%%%%%%%%%%%%%%%%%%%%%%%%%%%%%%%%%%%%%%%%%%%%%
